\begin{center}
\textbf{Posterior Integration and Volatility Inputs in WTI Asian Options}

\textbf{Integración posterior y volatilidad en opciones asiáticas WTI}
\end{center}

\section*{Abstract}

\textbf{Objective:} We examine the effect of historical-volatility posterior uncertainty on WTI Average Price Option values and compare it with the choice of volatility input. \textbf{Methods:} We combine a curvature-based approximation with 10 million synthetic return histories, posterior option pricing, and forward-in-time market comparisons. \textbf{Results:} The mean absolute posterior-integrated versus plug-in price difference is \$0.000745 per barrel in the October-2026 sample, while the mean width of the 95\% posterior theoretical-price interval is \$0.2733. In 2,381 out-of-sample observations, historical-volatility posterior integration has MAE 2.6187, versus 0.1075 for the previous-day APO implied-volatility smile. \textbf{Recommendation:} Benchmark historical inputs against strictly earlier option-implied information. \textbf{Limitations:} The sample covers few valuation dates and uses a one-factor constant-volatility model. \textbf{Originality:} We distinguish the point-price effect of posterior integration from posterior price dispersion and volatility-source effects in contract-consistent APO valuation. \textbf{Conclusion:} The volatility input matters more for observed pricing errors than posterior integration alone.

\textbf{Keywords:} Bayesian option pricing; Asian options; WTI crude oil; volatility uncertainty; implied volatility.

\section*{Resumen}

\textbf{Objetivo:} Evaluamos cómo la incertidumbre posterior de la volatilidad histórica afecta la valuación de opciones asiáticas sobre WTI y comparamos ese efecto con la selección de volatilidad. \textbf{Métodos:} Utilizamos una aproximación basada en la curvatura, 10 millones de muestras sintéticas, valuación posterior y comparaciones con información de opciones de fechas anteriores. \textbf{Resultados:} En la muestra de octubre de 2026, la diferencia absoluta media entre integración posterior y sustitución por la media posterior es de \$0.000745 por barril; la amplitud media del intervalo posterior de 95\% para precios teóricos es \$0.2733. En 2,381 observaciones fuera de muestra, el error absoluto medio con volatilidad histórica es 2.6187 frente a 0.1075 con volatilidad implícita del día anterior. \textbf{Recomendación:} Contrastar estimaciones históricas con información implícita estrictamente anterior. \textbf{Limitaciones:} Pocas fechas de valuación y un modelo de volatilidad constante con un factor. \textbf{Originalidad:} Distinguimos la corrección puntual por integración, la dispersión posterior de precios y el efecto de la fuente de volatilidad. \textbf{Conclusión:} La fuente de volatilidad domina las diferencias de error observadas.

\textbf{Palabras clave:} valuación bayesiana de opciones; opciones asiáticas; petróleo WTI; incertidumbre de volatilidad; volatilidad implícita.

\textbf{JEL:} C11; C15; C58; G13; G17; Q40.
